%%%
%%% All author customizations should go in this file
%%%

%%%---------------------------------
%%% DRAFT watermark and commentary control
%%%
%%% These two lines, if uncommented, will remove the watermark and margin comments
%%% to prepare your document for final formatting
\SetWatermarkText{}
% \setboolean{hidecomments}{true}

%%%---------------------------------
%%% Customize location of Model-Generated content, if using MDSA tools
%%%
%%% Uncomment this and provide the path to *model-generated* .tex and .svg content
%%% if you are placing somewhere other than the standard locations entered here...
%  \renewcommand{\genfiles}{GeneratedContent/}
%  \graphicspath{{../GeneratedContent/Images/}}


%%%---------------------------------
%%% Bibliography References
%%% See https://www.overleaf.com/learn/latex/Bibliography_management_in_LaTeX
%%% for guidance
%%%
%%% Location of the .bib file to provide your own bibliography entries
%%% If you have additional .bib files you may add them by copying the line and
%%% updating the filename. Note that you must include ".bib" at the end of the
%%% filename.
\addbibresource{specification.bib}

%%% OMG publication guidelines strongly suggest that references be broken out into
%%% normative, and non-normative lists.  By default, cited references will appear
%%% in the NON-normative bibliography. Add bibliographic citation references to 
%%% this list to have them appear in the normative bibliography instead.
% \addtocategory{normative}{MOF25, MOF251, UML251, XMI, XMLSchema}
\addtocategory{normative}{BMM, BPMN20, CQL, IEEE754, ISO8601, ISOEBNF, Java, ONNX, PMML,
  RFC3986, SCE10, UML241, XBASE, XML, XSD2, XDM, XFO}

%%% LaTeX will build your bibliography section from the references you cite in your
%%% document. If you wish to include certain 'default' references without needing to 
%%% explicitly find a place to add a \cite command, you may add them to the following
%%% list instead.
% \nocite{MOF25, MOF251, UML251, XMI, XMLSchema}
%%% DMN lists all its references in Clause 3, also those it does not cite.
\nocite{BMM, BPMN20, CQL, IEEE754, ISO8601, ISOEBNF, Java, ONNX, PMML, RFC3986, SCE10,
  UML241, XBASE, XML, XSD2, XDM, XFO, JSON, PRR, RIF, SBVR, SQL, XPath}


%%%---------------------------------
%%% Useful but completely optional packages
%%% Feel free to comment out the following sections if you are not using them.

%%% Source code listings, with useful defaults
%%% https://www.overleaf.com/learn/latex/Code_listing#Using_listings_to_highlight_code
\usepackage{listings}
\lstset{breaklines=true, tabsize=4, captionpos=b, basicstyle=\ttfamily, showstringspaces=false, keepspaces=true}

%%% Characters that pdfLaTeX's fonts have no glyph for, such as emoji, need a
%%% mapping, or the build stops with "Unicode character ... not set up for use
%%% with LaTeX" in text and "Invalid UTF-8 byte sequence" in code. Each character
%%% is declared twice: with \newunicodechar for text, and with the literate
%%% option of listings for code. Packages newunicodechar and twemojis are listed
%%% in user-pkgs.txt.
\usepackage{newunicodechar}
\usepackage{twemojis}
%%% U+1F40E horse, in DMN's FEEL string examples
\newunicodechar{🐎}{\twemoji{1f40e}}
\lstset{literate={🐎}{{\twemoji{1f40e}}}1}

%%% Subscripts and superscripts. The Markdown files write them as Unicode
%%% characters (e₁, eₙ, 10⁻¹), which GitHub shows; pdfLaTeX prints them with
%%% \textsubscript and \textsuperscript.
\newunicodechar{₀}{\textsubscript{0}}
\newunicodechar{₁}{\textsubscript{1}}
\newunicodechar{₂}{\textsubscript{2}}
\newunicodechar{₃}{\textsubscript{3}}
\newunicodechar{₄}{\textsubscript{4}}
\newunicodechar{₅}{\textsubscript{5}}
\newunicodechar{₆}{\textsubscript{6}}
\newunicodechar{₇}{\textsubscript{7}}
\newunicodechar{₈}{\textsubscript{8}}
\newunicodechar{₉}{\textsubscript{9}}
\newunicodechar{₋}{\textsubscript{-}}
\newunicodechar{ᵢ}{\textsubscript{i}}
\newunicodechar{ⱼ}{\textsubscript{j}}
\newunicodechar{ₖ}{\textsubscript{k}}
\newunicodechar{ₘ}{\textsubscript{m}}
\newunicodechar{ₙ}{\textsubscript{n}}
\newunicodechar{⁰}{\textsuperscript{0}}
\newunicodechar{⁴}{\textsuperscript{4}}
\newunicodechar{⁵}{\textsuperscript{5}}
\newunicodechar{⁶}{\textsuperscript{6}}
\newunicodechar{⁷}{\textsuperscript{7}}
\newunicodechar{⁸}{\textsuperscript{8}}
\newunicodechar{⁹}{\textsuperscript{9}}
\newunicodechar{⁻}{\textsuperscript{-}}
\newunicodechar{⁽}{\textsuperscript{(}}
\newunicodechar{⁾}{\textsuperscript{)}}
\newunicodechar{ˢ}{\textsuperscript{s}}

%%% Mathematical and other symbols in the text
\newunicodechar{→}{\ensuremath{\rightarrow}}
\newunicodechar{≤}{\ensuremath{\leq}}
\newunicodechar{≥}{\ensuremath{\geq}}
\newunicodechar{≡}{\ensuremath{\equiv}}
\newunicodechar{ε}{\ensuremath{\varepsilon}}
\newunicodechar{−}{\ensuremath{-}}
\newunicodechar{▪}{\rule{0.5ex}{0.5ex}}
%%% U+1F40 as in the Word source of DMN 1.7, where U+1F40E was meant
\newunicodechar{ὀ}{\`{o}}

%%% Support for tables that span across multiple pages
%%% https://www.overleaf.com/learn/latex/Tables#Multi-page_tables
\usepackage{longtable}

%%%---------------------------------
%%% User added packages
%%%
%%% You may add any additional changes you like below this line
